\chapter{Listados de referencia}
\label{chap:listados}

\lettrine{E}{n} este anexo se recogen íntegras las definiciones de los mensajes propios del sistema (sección~\ref{subsec:mensajes-propios}) y los ficheros que definen su configuración y despliegue (sección~\ref{sec:configuracion_despliegue}). Los tres listados que se emplean como ejemplo en el cuerpo de la memoria —el mensaje \texttt{CarLocation} y el fichero de lanzamiento y el \emph{compose} del nodo de captura— están en los listados~\ref{lst:msg-carlocation}, \ref{lst:camera-launch} y~\ref{lst:compose-camera}.

\section{Definición de los mensajes}
\label{sec:anexo-mensajes}

\begin{lstlisting}[caption={Mensajes que transportan la información del sistema.}, label={lst:msgs-informacion}]
# ===== FinishLine.msg =====
# Identificar que camara vio la linea de meta
string camara_id

LineSegment finish_line

# ===== SpeedCarril.msg =====
int32 pwm
# ri -> Forma de referirnos al carril
string carril
# Marca temporal para controlar cada cuanto
# nos llegan los mensajes del robot
builtin_interfaces/Time stamp

# ===== TimePerLap.msg =====
float64 lap_time
builtin_interfaces/Time stamp
int32 lap_number

# ===== CarControlTelemetry.msg =====
float64 dist_derrape
bool estado_derrapando
# Tiempo que pasa desde que recibimos la informacion
# del coche hasta que enviamos el pwm a RaceController
float64 pipeline_time
# Momento exacto en el que llego el paquete de la
# camara al CarController
builtin_interfaces/Time receive_msg_stamp
# Para saber si estamos detectando correctamente el coche
int32 roi_size
\end{lstlisting}

\begin{lstlisting}[caption={Mensajes auxiliares, empleados como piezas de construcción de los anteriores.}, label={lst:msgs-auxiliares}]
# ===== Point2D.msg =====
float64 x
float64 y

# ===== Stiker.msg =====
Point2D center
string color

# ===== BoundingRect.msg =====
int32 x
int32 y
int32 w
int32 h

# ===== LineSegment.msg =====
Point2D start
Point2D end
\end{lstlisting}

\section{Fichero de parámetros}
\label{sec:anexo-params}

\begin{lstlisting}[caption={Fichero de parámetros común del sistema (\texttt{config/params.yaml}).}, label={lst:params-yaml}]
/**:
  ros__parameters:
    modo_calibracion: True
    # Lista de coches del sistema. Genera un detector por coche en la camara,
    # un controlador por coche y las suscripciones del puente Arduino. Cada
    # nombre de esta lista tiene que tener su bloque en 'cars' de abajo.
    coches: ["car1", "car2"]

    # Configuracion POR COCHE. La clave (car1, car2, ...) es el identificador
    # del coche y coincide con 'coches' y con el namespace ROS (/car1/...).
    # Aqui vive TODO lo que distingue a un coche de otro:
    #   - stiker_front/back: colores de sus dos pegatinas. La camara construye
    #     un ColorDetector por coche con estos colores, y por eso dos coches
    #     con colores distintos se pueden seguir a la vez sin confundirse.
    #   - carril: carril fisico del Arduino al que esta conectado este coche
    #     ("r1"/"1" o "r2"/"2"; el puente parsea ambas formas). Sustituye al
    #     antiguo reparto por orden de la lista en el launch: aqui se pone el
    #     carril REAL del montaje, no el que toque por posicion.
    #   - modo_manual: false -> coche autonomo, lo arranca BrainLaunch (que
    #     ademas lanza el arduino_bridge). true -> lo conduce una persona, lo
    #     arranca ManualLaunch (sin puente). Cada launch filtra 'coches' por
    #     este flag, asi que para un mixto (uno manual + uno autonomo) se
    #     levantan los dos docker-compose a la vez y cada uno coge su parte.
    cars:
      car1:
        stiker_front: "azul"
        stiker_back: "verde"
        carril: "2"
        modo_manual: false
      car2:
        stiker_front: "verde"
        stiker_back: "azul"
        carril: "1"
        modo_manual: true

    #Parametros para la camara
    camera:
      # Controlar el tamano de la foto que sacamos
      # CAMERA_MODES = {
      #     0: (160, 120),
      #     1: (320, 240),
      #     2: (640, 480),
      #     3: (800, 600),
      #     4: (1280, 720),
      # }
      mode: 2
      device: '0'
      # Tamano del kernel que vamos a usar para generar la mascara
      mascara_kernel_size: 25
      # Parametro que controla si estamos en modo debug o no, en caso de
      # estar en debug se envia la imagen de lo que ve la camara
      debug: True

      # Parametros para el modulo de deteccion
      # (los colores de las pegatinas NO van aqui: son por coche, en 'cars')
      detection:
        # Color de la linea del meta del circuito
        finish_line_color: "naranja"
        # Tamano del kernel que usamos para la operacion de deteccion
        kernel_size: 5
        # Area minima que tiene que tener el stiker detectado para ser valido
        min_area: 40

    # Parametros para el controlador
    controller:
      # Velocidad minima a la que podemos ir
      minimum_speed: 65
      # Velocidad maxima que podemos alcanzar aplicando el algoritmo
      maximum_speed: 90
      distancia_nodos_trayectoria: 15.0
      umbral_control_trayectoria_correcta: 60.0
      umbral_derrape_peligro: 10.0
      umbral_derrape_seguro: 9.0
      latencia_min_nodos: 1
      latencia_max_nodos: 10

    # Parametros para el Arduino
    arduino:
      port: "/dev/ttyACM0"
      baudrate: 115200
      calibration_speed: 60
\end{lstlisting}

\section{Ficheros de lanzamiento}
\label{sec:anexo-launch}

\begin{lstlisting}[language=Python, caption={Lanzamiento de los controladores y del nodo de actuación (\texttt{launch/BrainLaunch.py}).}, label={lst:brain-launch}]
from launch import LaunchDescription
from launch_ros.actions import Node
from ament_index_python.packages import get_package_share_directory
import yaml
import os


def generate_launch_description():
    pkg_share = get_package_share_directory("image_processor_pkg")
    params_file = os.path.join(pkg_share, "config", "params.yaml")

    # 1. Cargamos la lista de coches y su configuracion por coche desde el YAML
    with open(params_file, "r") as f:
        config = yaml.safe_load(f)
        params = config["/**"]["ros__parameters"]
        coches = params["coches"]
        cars = params["cars"]

    ld = LaunchDescription()

    # 2. Un controlador por cada coche AUTONOMO (modo_manual: false). Los coches
    #    marcados como manuales los arranca ManualLaunch; para una carrera mixta
    #    (uno manual + uno autonomo) se levantan los dos docker-compose a la vez
    #    y cada launch coge su subconjunto del mismo params.yaml.
    hay_autonomo = False
    for car_name in coches:
        cfg = cars[car_name]
        if cfg.get("modo_manual", False):
            continue
        hay_autonomo = True

        ld.add_action(
            Node(
                package="image_processor_pkg",
                executable="CarControllerNode.py",
                # Nombre UNICO por coche: EstrategiaPerfil nombra sus logs con el
                # nombre del nodo (derrapesLog_<nodo>_<camara>.txt); con el mismo
                # nombre dos coches se pisarian los ficheros.
                name=f"CarController_{car_name}",
                namespace=car_name,
                parameters=[
                    params_file,
                    {
                        "car_name": car_name,
                        # Carril fisico del coche (cars.<coche>.carril), no por
                        # orden de la lista: el controlador lo mete en el mensaje
                        # SpeedCarril y el arduino_bridge lo usa para el carril.
                        "carril_asignado": str(cfg["carril"]),
                        "modo_manual": False,
                    },
                ],
                respawn=True,
                respawn_delay=2.0,
            )
        )

    # 3. El arduino_bridge solo se lanza si hay algun coche autonomo que mande
    #    PWM. Si todos los coches son manuales, no aporta nada y no se arranca.
    if hay_autonomo:
        ld.add_action(
            Node(
                package="image_processor_pkg",
                executable="RaceControllerNode.py",
                name="arduino_bridge",
                parameters=[params_file],
            )
        )

    return ld
\end{lstlisting}

\begin{lstlisting}[language=Python, caption={Lanzamiento en modo manual, sin nodo de actuación (\texttt{launch/ManualLaunch.py}). Se omite la cabecera de documentación del fichero.}, label={lst:manual-launch}]
from launch import LaunchDescription
from launch_ros.actions import Node
from ament_index_python.packages import get_package_share_directory
import yaml
import os


def generate_launch_description():
    pkg_share = get_package_share_directory("image_processor_pkg")
    params_file = os.path.join(pkg_share, "config", "params.yaml")

    # 1. Cargamos la lista de coches y su configuracion por coche desde el YAML
    with open(params_file, "r") as f:
        config = yaml.safe_load(f)
        params = config["/**"]["ros__parameters"]
        coches = params["coches"]
        cars = params["cars"]

    ld = LaunchDescription()

    # 2. Un controlador por coche MANUAL (modo_manual: true), igual que en
    #    BrainLaunch pero en manual. carril_asignado se mantiene aunque en
    #    manual no se use (publicar_velocidad sale antes): asi el nodo es el
    #    mismo en los dos modos y el mensaje SpeedCarril sigue llevandolo.
    for car_name in coches:
        cfg = cars[car_name]
        if not cfg.get("modo_manual", False):
            continue

        ld.add_action(
            Node(
                package="image_processor_pkg",
                executable="CarControllerNode.py",
                name=f"CarControllerManual_{car_name}",
                namespace=car_name,
                parameters=[
                    params_file,
                    {
                        "car_name": car_name,
                        "carril_asignado": str(cfg["carril"]),
                        "modo_manual": True,
                    },
                ],
                respawn=True,
                respawn_delay=2.0
            )
        )

    # 3. Aqui NO va el arduino_bridge: ver el punto 1 de la cabecera.

    return ld
\end{lstlisting}

\section{Ficheros de despliegue en contenedores}
\label{sec:anexo-compose}

\begin{lstlisting}[caption={Despliegue de los controladores y del nodo de actuación en la máquina conectada al Arduino (\texttt{docker-compose-brain.yml}).}, label={lst:compose-brain}]
services:
  cerebro:
    build: .
    network_mode: "host"
    privileged: true
    ipc: "host"
    pid: "host"
    devices:
      - "/dev/ttyACM0:/dev/ttyACM0"
    environment:
      - ROS_DOMAIN_ID=42
      - ROS_LOCALHOST_ONLY=0
      # Sin esto el contenedor va en UTC (ros:humble viene asi) y las marcas
      # de hora de los logs del algoritmo (derrapesLog_*.txt) salen
      # desfasadas respecto al reloj del PC y respecto al analisis, que corre
      # fuera del contenedor. Los stamps de los mensajes NO se ven afectados:
      # son epochs y las restas entre ellos dan igual el huso
      - TZ=Europe/Madrid
    volumes:
      # La zona horaria del host, ya compilada: glibc la lee aunque la imagen
      # no tuviese tzdata. Va junto a TZ porque cada mecanismo cubre un hueco
      # del otro (ver comentario de arriba)
      - /etc/localtime:/etc/localtime:ro
      - ./src/image_processor_pkg:/ros2_ws/src/image_processor_pkg
    command: >
      bash -c "source /opt/ros/humble/setup.bash &&
               colcon build --symlink-install &&
               source install/setup.bash &&
               ros2 launch image_processor_pkg BrainLaunch.py"
\end{lstlisting}

\begin{lstlisting}[caption={Despliegue en modo manual: sin acceso al Arduino y sin nodo de actuación (\texttt{docker-compose-manual.yml}).}, label={lst:compose-manual}]
# Despliegue del cerebro para una CARRERA MANUAL (la conduce una persona con
# el mando fisico del Scalextric). Es docker-compose-brain.yml con dos
# diferencias:
#
#   1. Sin el bloque `devices`: aqui el Arduino no se usa (queda fuera del
#      circuito), asi que el contenedor arranca aunque no este ni conectado.
#   2. Lanza ManualLaunch.py, que arranca los mismos controladores pero SIN
#      el arduino_bridge y con modo_manual=True.
#
# Uso:  docker compose -f docker-compose-manual.yml up
services:
  cerebro:
    build: .
    network_mode: "host"
    privileged: true
    ipc: "host"
    pid: "host"
    environment:
      - ROS_DOMAIN_ID=42
      - ROS_LOCALHOST_ONLY=0
      - TZ=Europe/Madrid
    volumes:
      - /etc/localtime:/etc/localtime:ro
      # Monta el paquete desde el host: sin esto los derrapesLog_*.txt que
      # escribe el algoritmo (en src/image_processor_pkg/logs_carrera) se
      # perderian al parar el contenedor, y son justo lo que hay que pasarle
      # al analisis para ver que habria hecho el algoritmo
      - ./src/image_processor_pkg:/ros2_ws/src/image_processor_pkg
    command: >
      bash -c "source /opt/ros/humble/setup.bash &&
               colcon build --symlink-install &&
               source install/setup.bash &&
               ros2 launch image_processor_pkg ManualLaunch.py"
\end{lstlisting}
